\begin{center}
\textbf{Abstract}
\end{center}

Demand planning teams optimize forecast accuracy; operations teams execute plans. These objectives diverge: a more accurate forecast often requires plan changes that exceed operational absorption capacity, creating coordination burden and execution gridlock. This paper reframes forecast-driven demand planning as feasibility-constrained model selection, evaluating forecast accuracy and execution feasibility as dual objectives.

The core contribution is a forecast-to-plan feasibility gate: candidate forecasts are converted into executable planning signals, measured for operational consequence (inventory exposure, plan stability, execution violations, governance burden), and ranked by feasibility-adjusted planning utility rather than accuracy alone. The gate is model-agnostic, enabling fair comparison of accuracy-first, ensemble, smoothing, greedy, and dynamic programming policies.

Empirically, the study draws on DataCo fulfillment records to calibrate data-informed execution-risk severity levels, alongside Favorita and M5 retail demand data. Across 123 evaluation groups, the accuracy-first deployable strategy differs from the operational-loss-optimal deployable strategy in 91.1\% of cases. The accuracy-first planning-execution gap rate is 15.4\% of periods, compared with 2.7\% for the best deployable operational-loss strategy. Ensemble methods offer low-disruption alternatives; dynamic programming and smoothing serve different governance and execution constraints, and strategy selection varies with planning hierarchy and demand sparsity.

The practical implication is direct: deployment decisions should be made at the forecast-to-plan layer, not the forecasting layer. Forecast accuracy remains important but should be treated as a diagnostic metric, separate from operational planning utility. Organizations can improve plan feasibility while sacrificing only modest accuracy, enabling faster, more confident model rollouts.

\begin{center}
\textbf{Keywords}
\end{center}

Demand planning, Forecast-to-plan conversion, Operational feasibility, Planning stability, Execution risk.

